\section{Fukaya-Theoretic Interpretation}
The coherent--constructible correspondence proved in the previous section was
formulated entirely in terms of coherent sheaves and constructible sheaves.
The proof used the FLTZ skeleton as a singular-support condition and avoided
any explicit use of Fukaya categories. Nevertheless, the constructible side
has a natural symplectic interpretation: constructible sheaves on
\(M_{\mathbb R}\) with prescribed singular support may be regarded as a
sheaf-theoretic model for a Fukaya category of the cotangent bundle
\(T^*M_{\mathbb R}\).

The purpose of this section is to explain this interpretation. We will not
construct the Fukaya category from first principles. Instead, we use the
microlocal sheaf--Fukaya correspondence as a guiding principle: a conic
Lagrangian subset of a cotangent bundle controls both singular supports of
constructible sheaves and asymptotic supports of Lagrangian branes. In our
case, the relevant conic Lagrangian is precisely the FLTZ skeleton
\(\Lambda_\Sigma\).

Thus the toric CCC may be viewed as part of a larger triangle of
correspondences:
\[
\operatorname{Perf}_T(X_\Sigma)
\simeq
Sh_{\Lambda_\Sigma}(M_{\mathbb R})
\simeq
\mathcal W(T^*M_{\mathbb R},\Lambda_\Sigma^\infty),
\]
where the right-hand side denotes a wrapped or partially wrapped Fukaya
category with stop determined by the contact boundary
\(\Lambda_\Sigma^\infty\) of the FLTZ skeleton. This equivalence should be
understood up to the usual conventions concerning compact objects, brane
structures, and possible opposite categories.

In this sense, the constructible side is not merely a combinatorial
replacement for coherent sheaves. It is also a microlocal model for the
Fukaya-theoretic mirror of the toric variety. The following subsections
explain how the singular-support condition becomes a Lagrangian support
condition, how microlocal sheaves are related to wrapped Fukaya categories,
and how the theta sheaves appearing in the CCC may be interpreted as
Lagrangian branes.


\subsection{From Singular Support to Lagrangian Support}

In Section 4, the FLTZ skeleton
\(\Lambda_\Sigma\subset T^*M_{\mathbb R}\) was introduced as the
singular-support condition defining the constructible side of the CCC. The
purpose of this subsection is to explain how the same subset is interpreted
on the Fukaya-theoretic side.

\begin{definition}[Cotangent symplectic structure]
Let \(Y\) be a smooth real manifold. The cotangent bundle \(T^*Y\) carries
its canonical Liouville one-form \(\lambda\). Its exterior derivative
\[
\omega=d\lambda
\]
is the canonical exact symplectic form on \(T^*Y\).
\end{definition}

\begin{definition}[Lagrangian submanifold]
Let \((X,\omega)\) be a symplectic manifold. A submanifold
\(L\subset X\) is called Lagrangian if \(\omega|_L=0\) and
\(\dim L=\frac{1}{2}\dim X\). In particular, a Lagrangian submanifold of
\(T^*Y\) has dimension \(\dim Y\).
\end{definition}

\begin{definition}[Lagrangian brane]
In Fukaya theory, one usually works not merely with Lagrangian
submanifolds, but with Lagrangians equipped with additional structures such
as orientations, gradings, spin structures, and local systems. Such an
enhanced Lagrangian is called a Lagrangian brane. In this section, the word
``brane'' will be used only in this sense.
\end{definition}

\begin{definition}[Boundary at infinity of a conic subset]
Let \(\Lambda\subset T^*Y\) be a conic subset, namely a subset invariant
under positive scaling in the cotangent fibers. Its boundary at infinity is
\[
\Lambda^\infty
:=
(\Lambda\setminus T_Y^*Y)/\mathbb R_{>0}
\subset S^*Y,
\]
where \(T_Y^*Y\) denotes the zero section and
\[
S^*Y:=(T^*Y\setminus T_Y^*Y)/\mathbb R_{>0}
\]
is the cosphere bundle.
\end{definition}

For \(Y=M_{\mathbb R}\), we have the identification
\(T^*M_{\mathbb R}\simeq M_{\mathbb R}\times N_{\mathbb R}\). Under this
identification, the fiberwise scaling action is the scaling action on the
\(N_{\mathbb R}\)-factor. By Proposition 4.2, the FLTZ skeleton
\(\Lambda_\Sigma\) is conic. Hence its boundary at infinity
\[
\Lambda_\Sigma^\infty
:=
(\Lambda_\Sigma\setminus T^*_{M_{\mathbb R}}M_{\mathbb R})/\mathbb R_{>0}
\subset S^*M_{\mathbb R}
\]
is well-defined.

On the sheaf-theoretic side, a constructible sheaf \(F\) has singular
support \(SS(F)\subset T^*M_{\mathbb R}\). This subset records the cotangent
directions in which local sections of \(F\) fail to propagate. Therefore the
condition
\[
SS(F)\subset \Lambda_\Sigma
\]
means that \(F\) is allowed to have singular behavior only in the
microlocal directions prescribed by the FLTZ skeleton.

On the Fukaya-theoretic side, the analogous restriction is imposed on
Lagrangian branes in \(T^*M_{\mathbb R}\). The subset
\(\Lambda_\Sigma^\infty\subset S^*M_{\mathbb R}\) records the directions at
infinity determined by the FLTZ skeleton. In wrapped Fukaya theory, such a
subset is used as a stop: it restricts the allowed wrapping behavior of
Lagrangian branes at infinity.

Thus the same geometric object has two interpretations. On the
constructible side, \(\Lambda_\Sigma\) is a singular-support constraint. On
the Fukaya side, \(\Lambda_\Sigma^\infty\) is an asymptotic Lagrangian
constraint. This is the basic bridge from constructible sheaves to Fukaya
categories; the precise categorical form of this bridge will be recalled in
the next subsection.

\subsection{Microlocal Sheaves and Wrapped Fukaya Categories}

We now recall the external theorem relating constructible sheaves to Fukaya
categories. The terminology concerning conic subsets, their boundary at
infinity, and stops was fixed in the previous subsection. This theorem is
not proved in this paper; it is used only to interpret the constructible side
of the CCC as a Fukaya-theoretic category.

Let \(Y\) be a real analytic manifold, and let
\(\Lambda\subset T^*Y\) be a closed conic subanalytic Lagrangian subset.
We write \(Sh_\Lambda(Y;k)^c\) for the compact objects in the dg category of
constructible \(k\)-sheaves on \(Y\) whose singular support is contained in
\(\Lambda\). We write
\(\mathcal W(T^*Y,\Lambda^\infty;k)\) for the partially wrapped Fukaya
category of \(T^*Y\) stopped at \(\Lambda^\infty\), with the brane and
grading conventions fixed.

\begin{theorem}[Microlocal sheaf--Fukaya correspondence, quoted form]
With the notation above, there is an equivalence of dg categories
\[
Sh_\Lambda(Y;k)^c
\simeq
\mathcal W(T^*Y,\Lambda^\infty;k).
\]
\end{theorem}

We use this theorem as an external input. Depending on the convention for
Fukaya morphisms and for standard versus costandard sheaves, the right-hand
side may have to be replaced by its opposite category. This convention does
not affect the role of the theorem here.

Applying the theorem to \(Y=M_{\mathbb R}\) and
\(\Lambda=\Lambda_\Sigma\), we obtain
\[
Sh_{\Lambda_\Sigma}(M_{\mathbb R};k)^c
\simeq
\mathcal W(T^*M_{\mathbb R},\Lambda_\Sigma^\infty;k).
\]
Thus the constructible side of the toric CCC may be viewed as a microlocal
sheaf model for the stopped wrapped Fukaya category determined by the FLTZ
skeleton.

\subsection{Theta Sheaves as Branes}

We now apply the microlocal sheaf--Fukaya correspondence from the previous
subsection to the theta sheaves appearing in the CCC. Let
\[
\Phi_\Sigma:
Sh_{\Lambda_\Sigma}(M_{\mathbb R};k)^c
\xrightarrow{\sim}
\mathcal W(T^*M_{\mathbb R},\Lambda_\Sigma^\infty;k)
\]
be an equivalence as in Theorem 6.1. We use this equivalence to transport
the constructible theta objects to the Fukaya side.

\begin{lemma}
For every \(\sigma\in\Sigma\) and \(\bar\chi\in M_\sigma\), the theta sheaf
\(\Theta(\sigma,\bar\chi)\) has singular support contained in
\(\Lambda_\Sigma\). Hence it defines an object of
\(Sh_{\Lambda_\Sigma}(M_{\mathbb R};k)\).
\end{lemma}

\begin{proof}
By definition, \(\Theta(\sigma,\bar\chi)\) is the costandard sheaf supported
on the closed convex polyhedron
\(P_{\sigma,\bar\chi}=\bar\chi+\sigma^\vee\). The singular support of such a
costandard polyhedral sheaf is controlled by the conormal cones to the faces
of \(P_{\sigma,\bar\chi}\). These conormal directions are precisely the
directions appearing in the FLTZ skeleton. Therefore
\(SS(\Theta(\sigma,\bar\chi))\subset\Lambda_\Sigma\).
\end{proof}

\begin{definition}
The elementary theta brane associated to \((\sigma,\bar\chi)\) is the Fukaya
object
\[
L_{\sigma,\bar\chi}
:=
\Phi_\Sigma\bigl(\Theta(\sigma,\bar\chi)\bigr)
\in
\mathcal W(T^*M_{\mathbb R},\Lambda_\Sigma^\infty;k).
\]
\end{definition}

Thus \(L_{\sigma,\bar\chi}\) is not introduced by an independent Floer-theoretic
construction in this paper. It is defined as the Fukaya object corresponding
to the theta sheaf under the microlocal sheaf--Fukaya equivalence. The
polyhedron \(P_{\sigma,\bar\chi}\) records the base support of the sheaf,
while the conormal geometry of its faces records the Lagrangian directions
controlled by the FLTZ skeleton.

\begin{proposition}[Morphisms between elementary theta branes]
Let \(\sigma,\tau\in\Sigma\), \(\bar\chi\in M_\sigma\), and
\(\bar\psi\in M_\tau\). Then
\[
\operatorname{RHom}_{\mathcal W}
\bigl(L_{\sigma,\bar\chi},L_{\tau,\bar\psi}\bigr)
\simeq
\operatorname{RHom}
\bigl(\Theta(\sigma,\bar\chi),\Theta(\tau,\bar\psi)\bigr).
\]
Consequently,
\[
\operatorname{RHom}_{\mathcal W}
\bigl(L_{\sigma,\bar\chi},L_{\tau,\bar\psi}\bigr)
\simeq
\begin{cases}
k, & P_{\sigma,\bar\chi}\subset P_{\tau,\bar\psi},\\
0, & \text{otherwise}.
\end{cases}
\]
In the nonzero case, the complex is concentrated in degree zero.
\end{proposition}

\begin{proof}
The first quasi-isomorphism follows from the fact that
\(\Phi_\Sigma\) is an equivalence of dg categories. The explicit computation
then follows from the Hom-complex calculation for constructible theta
sheaves proved in Section 5.2.
\end{proof}

Thus the theta sheaves become elementary Fukaya branes under the microlocal
sheaf--Fukaya correspondence. In the sheaf-theoretic proof of the CCC, the
objects \(\Theta(\sigma,\bar\chi)\) generate the constructible side. In the
Fukaya-theoretic interpretation, the corresponding objects
\(L_{\sigma,\bar\chi}\) form the elementary branes whose asymptotic behavior
is governed by \(\Lambda_\Sigma^\infty\).
\subsection{The Fukaya-Theoretic Form of the CCC}
We now combine the toric CCC with the microlocal sheaf--Fukaya
correspondence recalled in Section 6.2. The CCC proved in Section 5 gives an
equivalence between the equivariant coherent category of the toric variety
and the constructible category with singular support in the FLTZ skeleton.
The microlocal sheaf--Fukaya correspondence identifies this constructible
category with a stopped wrapped Fukaya category of the cotangent bundle.

\begin{corollary}[Fukaya-theoretic form of the toric CCC]
Let \(X_\Sigma\) be a smooth toric variety, and let
\(\Lambda_\Sigma\subset T^*M_{\mathbb R}\) be the FLTZ skeleton. After fixing
the Fukaya-theoretic brane and grading conventions, there is an equivalence
of dg categories
\[
\operatorname{Perf}_T(X_\Sigma)
\simeq
\mathcal W(T^*M_{\mathbb R},\Lambda_\Sigma^\infty;k).
\]
Equivalently, one has a chain of equivalences
\[
\operatorname{Perf}_T(X_\Sigma)
\simeq
Sh_{\Lambda_\Sigma}(M_{\mathbb R};k)^c
\simeq
\mathcal W(T^*M_{\mathbb R},\Lambda_\Sigma^\infty;k).
\]
\end{corollary}

\begin{proof}
The first equivalence in the chain is the compact form of the toric
coherent--constructible correspondence proved in Section 5. The second
equivalence is the microlocal sheaf--Fukaya correspondence from Theorem 6.1,
applied to \(Y=M_{\mathbb R}\) and
\(\Lambda=\Lambda_\Sigma\). Composing these equivalences gives the desired
Fukaya-theoretic form of the CCC.
\end{proof}

Thus the equivariant coherent category of the toric variety can be interpreted
as a stopped wrapped Fukaya category of the cotangent bundle
\(T^*M_{\mathbb R}\). Under this interpretation, the FLTZ skeleton plays the
role of the stop at infinity. On the constructible side, it appears as the
singular-support condition
\[
SS(F)\subset\Lambda_\Sigma.
\]
On the Fukaya side, its boundary at infinity
\(\Lambda_\Sigma^\infty\subset S^*M_{\mathbb R}\) controls the allowed
asymptotic behavior of Lagrangian branes.

The theta objects also admit a Fukaya-theoretic interpretation. The coherent
theta object \(\Theta'(\sigma,\bar\chi)\) corresponds under the CCC to the
constructible theta sheaf \(\Theta(\sigma,\bar\chi)\), and under the
microlocal sheaf--Fukaya correspondence this object corresponds to the
elementary theta brane \(L_{\sigma,\bar\chi}\). Therefore the correspondence
on generators can be summarized as
\[
\Theta'(\sigma,\bar\chi)
\longleftrightarrow
\Theta(\sigma,\bar\chi)
\longleftrightarrow
L_{\sigma,\bar\chi}.
\]
In this way, the Fukaya-theoretic form of the CCC identifies coherent theta
objects with elementary Lagrangian branes governed by the same polyhedral
data.

This formulation explains the role of the symplectic material from the
beginning of the paper. The moment map and moment polytope describe how
convex geometry arises naturally from toric symplectic geometry, while the
FLTZ skeleton gives the corresponding microlocal Lagrangian support condition
inside \(T^*M_{\mathbb R}\). The sheaf-theoretic CCC can therefore be viewed
as a bridge between equivariant coherent sheaves on \(X_\Sigma\) and Fukaya
objects controlled by the FLTZ skeleton.